\section{Avant l'électricité (avant 1794)}\label{sec:era0}
% =============================================================================

\subsection{L'écriture comptable et les supports de mémoire}\label{subsec:ecriture}

Les techniques de l'information ne commencent pas avec l'électricité. Elles commencent avec l'argile. Dans le sud de la Mésopotamie, entre l'Euphrate et le Tigre, l'État sumérien qui se forme au troisième millénaire avant notre ère développe un système d'écriture cunéiforme dont la fonction première n'est pas littéraire mais comptable~: enregistrer les stocks de grain, les surfaces irriguées, les journées de travail, les rations distribuées \parencite{nissen_archaic_1993}. Comme l'ont établi \textcite{nissen_archaic_1993} et \textcite{englund_hard_1991}, ce système informationnel naît des besoins d'une économie redistributive et donne à l'administration un instrument de contrôle sur des projets hydrauliques de grande échelle~--- canaux, digues, bassins de rétention. \dimmark{D7}{?}%
L'écriture rend l'irrigation gouvernable, et l'irrigation rend l'État viable. \dimmark{0-TER}{?}%
Mais les mêmes canaux, gérés sur plusieurs siècles, provoquent la salinisation des sols du sud mésopotamien \parencite{dickson_circumscription_1987}~: l'un des premiers épisodes documentés de destruction environnementale à l'échelle d'un écosystème est indissociable de la première technologie de l'information.

Parallèlement à cette trajectoire du stockage, le calcul a sa propre histoire. L'abaque apparaît sous des formes indépendantes dans des civilisations qui ne communiquent pas entre elles~: jetons en argile à Sumer ($\sim$2500~av.~J.-C.), tables à rainures en Grèce et à Rome, bouliers en Chine (dynastie Han) et au Japon (\emph{soroban}). \dimmark{D5}{?}%
Cette convergence suggère que la demande de calcul naît partout où la complexité économique~--- commerce, fiscalité, gestion des stocks~--- dépasse la capacité du calcul mental. Le système de numération positionnelle sexagésimale des Babyloniens, transmis par l'astronomie arabe puis repris par les astronomes européens, est encore celui des heures, des minutes et des degrés~: \dimmark{D10}{?}%
c'est un cas précoce de verrouillage de standard, où le coût de basculement maintient une convention en place sur quatre millénaires. Le mécanisme d'Anticythère ($\sim$200~av.~J.-C.), retrouvé dans une épave en 1901, est un calculateur astronomique à engrenages capable de prédire les éclipses et les positions planétaires \parencite{freeth_decoding_2006}. Mais il n'aura pas de descendance~: en l'absence d'une demande économique ou administrative suffisante pour justifier sa reproduction, le calcul mécanique sera réinventé au \textsc{xvii}\ieme{} siècle, pas hérité de l'Antiquité.

\subsection{Le transport de l'information}\label{subsec:transport}

La transmission d'information à distance repose sur le porteur humain et animal. Le \emph{cursus publicus} romain, organisé par Auguste, relie Rome aux provinces par des relais espacés de dix à quinze miles~; sa vitesse moyenne ne dépasse pas cinquante miles par jour \parencite{kolb_transport_2000}. Mille ans plus tard, la poste de la famille Taxis, à partir de 1490, couvre les Pays-Bas, l'Italie et l'Empire. La vitesse maximale de transmission de l'information reste celle du cheval au galop~--- environ trente kilomètres par jour en moyenne, relais compris. Le message ne voyage pas plus vite que le messager. \dimmark{D6}{?}%
La portée de la coordination économique~--- marchés, administrations, empires~--- était ainsi bornée par la vitesse du transport physique\footnote{Le principe est posé dès 1776 par Adam Smith~: la division du travail est limitée par l'étendue du marché (\emph{Wealth of Nations}, livre~I, chap.~3). La vitesse de transmission de l'information est l'une des contraintes sur cette étendue.}.

\subsection{L'imprimerie et la reproduction du savoir}\label{subsec:imprimerie}

La transformation du stockage intervient avec l'imprimerie à caractères mobiles de Gutenberg ($\sim$1440), \dimmark{D2}{?}%
qui réduit le coût de reproduction du texte de plusieurs ordres de grandeur\footnote{\textcite{eisenstein_printing_1979} a montré que cette chute du coût marginal de reproduction transforma non seulement la diffusion du savoir mais aussi les conditions de l'activité scientifique elle-même, en rendant possible la comparaison systématique des sources et la cumulativité des résultats.} sans modifier la nature du support~: le papier, introduit en Europe depuis la Chine via le monde arabe au \textsc{xii}\ieme{} siècle, reste un support inerte qui ne consomme pas d'énergie pour exister. La diffusion du \emph{Liber Abaci} de Fibonacci (1202), qui introduit la numération indo-arabe en Europe, et celle des tables de logarithmes de Napier (1614), sont des effets de cette infrastructure de stockage. L'imprimerie est aussi, comme le soutient \textcite{mokyr_gifts_2002}, l'infrastructure de la \og \emph{knowledge economy}\fg{} qui rend possible la révolution industrielle.

\subsection{Pascal, Leibniz et les premières machines à calculer}\label{subsec:pascaline}

Le calcul mécanique apparaît au \textsc{xvii}\ieme{} siècle avec deux prototypes sans descendance industrielle. La Pascaline (1642) de Blaise Pascal, conçue pour aider son père, commissaire aux impôts en Haute-Normandie, exécute l'addition et la soustraction par un mécanisme de roues à dix dents et de report de retenue. Le \emph{Stepped Reckoner} de Leibniz (1673) étend le mécanisme aux quatre opérations en introduisant le tambour cannelé (\emph{Leibniz wheel}), un cylindre à neuf dents de longueur croissante. Ni Pascal ni Leibniz ne résoudront le problème de la production en série~: la Pascaline sera fabriquée à une cinquantaine d'exemplaires, le \emph{Stepped Reckoner} restera à l'état de prototype. \dimmark{D5}{?}%
Comme pour le mécanisme d'Anticythère, la capacité technique précède la demande~: le calcul mécanique ne passera à l'échelle industrielle qu'avec l'Arithmometer de Thomas de Colmar, breveté en 1820, lorsque les compagnies d'assurance et les administrations fiscales fourniront un marché suffisant.

\subsection{Prony, le Nautical Almanac et la division du travail intellectuel}\label{subsec:prony}

Le calcul humain organisé, en revanche, précède de loin la machine et détermine les conditions de son apparition. En France, Gaspard de Prony organise en 1791 le calcul de tables logarithmiques et trigonométriques en mobilisant soixante-dix à quatre-vingts personnes réparties en trois niveaux hiérarchiques~: les mathématiciens qui conçoivent les formules, les calculateurs qualifiés qui les décomposent en opérations élémentaires, les calculateurs non qualifiés qui exécutent ces opérations. \dimmark{D4}{?}%
Babbage, qui visite le projet, en tirera dans \emph{On the Economy of Machinery and Manufactures} (1832) la conclusion que la division du travail de Smith s'applique aussi au traitement de l'information. En Angleterre, le Nautical Almanac Office, créé en 1767 par Nevil Maskelyne, emploie en permanence des \emph{computers}~--- le mot désigne alors une personne, pas une machine~--- pour calculer les éphémérides astronomiques nécessaires à la navigation\footnote{Les éphémérides ne sont pas un bien de consommation~: elles sont un intrant de la flotte marchande britannique, et leur production est financée par l'État pour des raisons de puissance commerciale et militaire. La demande de calcul est ici dérivée de la demande de commerce maritime.}.

\subsection{Le métier Jacquard et le contrôle par l'information stockée}\label{subsec:jacquard}

Le contrôle automatisé par l'information apparaît sous deux formes distinctes à la fin du \textsc{xviii}\ieme{} siècle. Le régulateur centrifuge de Watt (1788), qui coupe l'alimentation en vapeur quand la vitesse du moteur dépasse un seuil, est un mécanisme de rétroaction physique~--- du \emph{feedback} mécanique, pas du contrôle informationnel. Le métier Jacquard (1804), en revanche, est le premier cas de contrôle d'un processus physique par une séquence d'information stockée~: les cartes perforées encodent le motif du tissu, et le métier l'exécute sans intervention humaine. \dimmark{D3/D5}{?}%
L'innovation répond à une demande précise~: la complexité croissante des motifs exigés par le marché de la soie lyonnaise dépassait la capacité de coordination manuelle des tireurs de lacs\footnote{Avant le métier Jacquard, la réalisation d'un motif complexe nécessitait un \emph{tireur de lacs}, ouvrier spécialisé qui actionnait manuellement les fils de chaîne selon les instructions du dessinateur. Le métier Jacquard substituait du capital (les cartes, le mécanisme) à ce travail qualifié.}.

Ces trajectoires~--- transmettre, calculer, stocker, contrôler~--- coexistent sans converger. Les gens qui calculent les tables de Prony n'ont aucun rapport avec le réseau postal des Taxis, ni avec les imprimeurs de Gutenberg, ni avec les tisserands lyonnais qui travaillent sur les métiers Jacquard. L'histoire longue que cette section a esquissée n'est pas un prologue décoratif~: elle pose les conditions initiales à partir desquelles le télégraphe optique de Chappe (1794), puis le télégraphe électrique de Morse et Cooke \& Wheatstone (1837), vont pour la première fois mécaniser l'une de ces quatre fonctions~--- la transmission~--- et ouvrir la trajectoire que le reste du chapitre va parcourir.




% =============================================================================